% !TEX encoding = UTF-8 Unicode
% !BIB TS-program = biber
%
% This file is MIT-Thesis.tex, a LaTeX template for formatting an MIT thesis with the mitthesis class.
%
% Version: 1.24, 2026/07/27
%
% Author: John H. Lienhard, copyright 2026. Reuse under the MIT license: https://ctan.org/license/mit 
%
% DOCUMENTATION is here: https://ctan.org/pkg/mitthesis
%
% This package requires LaTeX formats dated 2022‑06‑01 or later


%% Don't remove the \DocumentMetadata command 
\DocumentMetadata
{
	lang		= en-US,          % default language
%
%	pdfversion  = 1.7,
%	pdfstandard = a-2b,
%
%	pdfversion  = 2.0,            % default version
%
	tagging		= on,             % For tagging, run lualatex with TeX Live 2026 or later.
					              % Tagging is necessary for accessibility and html rendering
					              % Do not include appendixb.tex when tagging is on (packages in that appendix are not compatible).
	pdfstandard = { ua-2, a-4f }, % requires pdfversion = 2.0 and TeX Live 2026 or later, with tagging on
	tagging-setup={math/setup=mathml-SE}, % math structural element tagging, needed for accessablity
}

%%%%%%%%%  Documentclass and options  %%%%%%%%%%%%%%%%%%%%%%%%%%%%%%%%%%%%%

\documentclass[twoside,fontset=termes-stix2]{mitthesis}
% 						fontset=libertinus, fontset=lmodern, fontset=lucida, 
%						fontset=fira-newtxsf, fontset=newtx, fontset=newtx-sans-text,  
%						fontset=heros-stix2,  fontset=stix2, fontset=termes-stix2, fontset=termes
%
% option [twoside]		gives facing-page behavior for printing; omitting twoside will eliminate even-numbered blank pages
%						(use mydesign to change page margins, e.g., for binding).
%
% option [lineno]	 	provides line numbers, as for editing
%
% option [fontset=xxx]  is a keyvalue which can be:
%					 	for pdftex or lualatex engine:  defaultfonts, libertinus, lmodern, lucida
%					 	for pdftex only: 				fira-newtxsf, newtx, newtx-sans-text
%						for lualatex only:				heros-stix2, stix2, termes, termes-stix2
%					 	if no key value is given, fonts default to CMR (pdftex) or LMR (unicode), i.e., "the LaTeX font".
%					 	You can edit the fontset files or you can write your own, myfonts.tex, and do [fontset=myfonts].
%						If you are using multiple languages, load the babel package in your fontset file, before the fonts.
%
% option [mydesign] 	loads packages for color, title and list formats, margins, or captions. 
%	or [mydesign=file]	You can edit mydesign.tex to change defaults. Other design files can be loaded as
%						key values, [mydesign=filename], where filename.tex should be in your working directory.
%						Two example files are provided: mydesign_redsans_headings.tex and mydesign_libertinus_headings.
%						See documentation for details and a description of the appropriate fontsets to use with each.
%
%					    As distributed, mydesign.tex will set caption labels in bold face (e.g., "Figure 1.1:" in bold)
%


%%%%%%%%  Packages used in the sample chapters (not otherwise required) %%%

%% Package for code listing in Appendix A.
\usepackage{listings}%   documentation is here https://ctan.org/pkg/listings

%% Set chemical formulas nicely
\usepackage[version=4]{mhchem}%   documentation at https://ctan.org/pkg/mhchem

%% Latin filler used in Chapter 1, with a test for package version date (https://ctan.org/pkg/lipsum)
\usepackage{lipsum}
\IfPackageAtLeastTF{lipsum}{2021/09/20}{\setlipsum{auto-lang=false}}{} % to avoid a hyphenation warning


%%%%%%%%%  Other optional packages used sample chapters  %%%%%%%%%%%%%%%%%%

\usepackage{microtype}% package for improved typography (https://ctan.org/pkg/microtype)

\usepackage{multicol}%  load if using the two-column nomenclature, \begin{nomenclature*}


%% Table related packages  

\usepackage{booktabs}% publication quality tables (https://ctan.org/pkg/booktabs)

\usepackage{array}%    additional options for column formats (https://ctan.org/pkg/array)

\usepackage{dcolumn}%  for alignment of numbers on the decimal place (https://ctan.org/pkg/dcolumn) 
  \newcolumntype{d}[1]{D{.}{.}{#1}}% use with dcolumn package. Note: dcolumns are set in math mode.
                                   % syntax: d{x.y} where x is max number of digits before decimal, y is max number after

%\usepackage{tabularx}% adjustable-width columns in tabular (https://ctan.org/pkg/tabularx)

\usepackage{longtable}% typeset multipage table, see Appendix A (https://ctan.org/pkg/longtable)
  \setlength{\LTcapwidth}{\linewidth}% width of caption over longtables.  Adjust as desired.

\usepackage[figuresright]{rotating}% to get sideways tables/figures, see Appendix A
								   % (https://ctan.org/pkg/rotating)
									

%%%%%%%%%  Graphics path (to figure files)  %%%%%%%%%%%%%%%%%%%%%%%%%%%%%%%%

%% Can set graphicspath to point to specific directories containing figures (the current directory is searched automatically)
%% For instance, to search a subdirectory of the current directory called "figures" and a parallel directory called "art", set:

% \graphicspath{ {figures/} {../art/} }% For details see: https://latexref.xyz/dev/latex2e.html#g_t_005cgraphicspath


%%%%%%%%%  Representative set-up for biblatex  %%%%%%%%%%%%%%%%%%%%%%%%%%%%%

%% biblatex is very powerful, and you can customize most aspects the reference list and citations to suit your needs.
%%   documentation is here: https://ctan.org/pkg/biblatex
%%   cheat sheet is here:   https://tug.ctan.org/info/biblatex-cheatsheet/biblatex-cheatsheet.pdf

%% Numerical citations of references
\usepackage[style=ext-numeric-comp,giveninits=true,maxbibnames=10,sorting=none,language=american]{biblatex}

% and make some changes to that style
	\renewcommand{\multicitedelim}{\addcomma}% removes space between consecutive cites to give [6,7], not [6, 7]
    \AtEveryBibitem{%
      \ifentrytype{article}{%
        \renewbibmacro{in:}{}% Removes "In:" for articles
		\renewcommand*{\jourvoldelim}{\addcomma\space}   % put comma after journal title
		\renewcommand*{\volnumdatedelim}{\addcomma\space}% put comma after volume info
		\renewbibmacro*{issue+date}{% Print date without parentheses around it
		\iffieldundef{issue}%
        	{}%
        	{\printfield{issue}%
         	\setunit*{\addcomma\space}
        	}%
			\printdate
		}
      }{}
    }
    %
	\renewbibmacro*{volume+number+eid}{%
        \printtext[bold]{\printfield{volume}}% bold volume (issue number) , eid
        \iffieldundef{number}
          {} % do nothing
          {\printtext[parens]{\printfield{number}}} % print number in parentheses
        \setunit{\bibeidpunct}%
        \printfield{eid}
    }

%% IEEE style citations and references
% \usepackage[style=ieee,maxbibnames=10,sorting=none]{biblatex}% style=ext-numeric-comp,articlein=false,giveninits=true
%	 \DefineBibliographyStrings{english}{url= \textsc{url} ,  }% replaces the IEEE default "[Online]. Available" by "URL"

%% author-year style citations and references 
%% use \parencite, not \cite, when you want "(Author, year)"
%% The sample files are not set up to include parentheses.
% \usepackage[style=authoryear, maxbibnames=10]{biblatex} 


\addbibresource{mitthesis-sample.bib}%% <== change to YOUR bib file <=============== CHANGE !!!!!


%% to avoid split urls and stretched white space, you can make the bibliography ragged-right by uncommenting this:
%\AddToHook{cmd/bibsetup/after}{\raggedright}

%% To ensure citations are set, run Latex --> biblatex --> Latex again (unless your installation does this automatically)


%%%%%%%%%%  Option for "double spacing" %%%%%%%%%%%%%%%%%%%%%%%%%%%%%%%%%%%%

%% Back in the typewriter era, double spaced lines were convenient for editing with a pencil. 
%% In typography, the separation between lines is called "leading", and it is usually set in 
%% proportion to the font size (i.e., when the font is loaded).  If you really feel the need 
%% to change the line separation, the most attractive results will be obtained by changing the
%% leading in proportion to the the current font size, rather than just doubling the space.

%% The setspace package provides a tool for changing line separation. Use these two commands here:
%
% \usepackage{setspace}%  documentation at https://ctan.org/pkg/setspace
% \setstretch{1.1}% you can choose some other value for the stretch of space between lines
%
%% Use one or more of the these commands *AFTER* the frontmatter (title page, abstract, acknowledgements, biosketch)
%
% \onehalfspacing
% \doublespacing
% \singlespacing  % will turn these effects off (you can use these anywhere in the document)

%% The best result is usually to stay with the default leading (or to learn about latex font settings).


%%%%%%%%%%%%  Math operators %%%%%%%%%%%%%%%%%%%%%%%%%%%%%%%%%%%%%%%%%%%%%%%%%%%%

% These commands declare two math operators, \erf{..} and \erfc{..} using the mathtools package.
% note: * form, \erf*{..}, produces automatic delimiter scaling; 
%       optional argument sets size manually, e.g. \erf[\bigg]{..}.
% See Table 1.1 in Chapter 1

\DeclareMathOperator{\Erf}{erf}
\DeclareMathOperator{\Erfc}{erfc}
\DeclarePairedDelimiterXPP\erf[1]{\Erf\mkern1mu}(){}{#1}% increase to 2mu with stix2 font
\DeclarePairedDelimiterXPP\erfc[1]{\Erfc\mkern1mu}(){}{#1}


%%%%%%%%%%%  Metadata (edit these to fit YOU)  %%%%%%%%%%%%%%%%%%%%%%%%%%%%%%%%%%%%%%%%%%%%%%%%%%%

% Most of the document metadata is created automatically. 
% The following items should be changed to match your work. <======================= !!!!!!!!!!!!!!!!!!

\hypersetup{%
	pdfsubject={Template for writing MIT theses with the mitthesis class},
	% Change this to briefly state topic of your thesis 
% 
	pdfkeywords={Massachusetts Institute of Technology, MIT},
	% Add keywords that will help search engines and libraries to find your work.
	% Includes the name[s] of the author[s] 
	% (If you used \DocumentMetadata at line 14, you can just put "\CopyrightAuthor," for the names.)
%
	pdfurl={},
	% If you have a url for the thesis, put it here. Otherwise delete this.
	% (MIT Libraries will put your thesis in DSPACE with a persistent url after you submit it.)
%	
	pdfcontactemail={},
	% You can put a [permanent] email address into the metadata, if you like.
	% Otherwise delete this.
%
	pdfauthortitle={},
	% If you have a title (e.g., Prof. Dr.-Ing.), you can include it here.
}

%%%%%%%%%%%%%%  End preamble %%%%%%%%%%%%%%%%%%%%%%%%%%%%%%%%%%%%%%%%%%%%%%%%%%%%%%%%%%%%%%%%%%%%%
%%%%%%%%%%%%%%%%%%%%%%%%%%%%%%%%%%%%%%%%%%%%%%%%%%%%%%%%%%%%%%%%%%%%%%%%%%%%%%%%%%%%%%%%%%%%%%%%%%

\begin{document}
%%% edit the following commands to match your thesis %%%%%%%%%%

\title{The Atomic Theory as Applied To Gases, with Some Experiments on the Viscosity of Air}

% \Author{Author full name}{Author department}[Author's first PREVIOUS degree][Author's second PREVIOUS degree][...
% Note that third, fourth, fifth, and sixth arguments are optional [] and may be omitted

\Author{Silas W. Holman}{Department of Physics}
  
% \Author{Silas W. Holman}{Department of Physics \\ & Department of Mathematics} % if author is in more than one department
% \Author{Luisa Hernández}{Department of Research}[B.S. Mechanical Engineering, UCLA, 2018][M.S. Stellar Interiors, Vulcan Science Academy, 2020] 
% \Author{Thurston Howell III}{Department of Economics}[M.B.A. Ferengi School of Management, 2022]
%
% note on these and other names: most of the names are made up, but Silas Holman was a physics professor at MIT in the 19th century.


% Use once for each degree fulfilled by thesis
% For two degrees from one department, leave the department argument blank for the second degree {}.
% For one degree from two departments, leave the degree argument blank for the second department {}.

\Degree{Bachelor of Science in Physics}{Department of Physics}

% \Degree{Master of Science in Physics}{}
% \Degree{Bachelor of Science in Mechanical Engineering}{Department of Mechanical Engineering}
%
% A very long degree name can be split into multiple lines as shown:
% \Degree{Doctor of Philosophy \\ in \\ Electrical Engineering and Computer Science}{Department of Electrical Engineering and Computer Science}


% If there is more than one supervisor, use the \Supervisor command for each.
% NB: optional department field added with v1.21 in late 2025. It is needed ONLY on the Thesis Committee page.

\Supervisor{Edward C. Pickering}{Professor of Physics}[Department of Physics]

% Separate multiple titles or departments with "\\ &".  An "\&" used by itself will be converted to "and". 
% \Supervisor{Secunda Castor}{Professor of Research \\ & Professor of Knowledge \\ & Professor of Reason}[Department of Research]
% \Supervisor{Quintus Castor}{Professor of Hydraulics}[Department of Civil \& Environmental Engineering \\ & Department of All Faculty Who Are Not Members of a Department] 


% Professor who formally accepts theses for your department (e.g., Professor Sméagol, Graduate Officer, Department of Powerful Rings)
% If more than one department gives the degree, use once for each department

\Acceptor{Primus Castor}{Professor of Log Dams}{Undergraduate Officer, Department of Physics}{}

% Separate multiple titles or departments with "\\ &".  An "\&" used by itself will be converted to "and". 
% \Acceptor{Tertius Castor}{Professor of Log Dams \\ & Professor of Hydraulics}{Graduate Officer, Department of Research}
%
%  If you need to reduce vertical space on title page, put the acceptor title in the second argument and leave the third blank, {}.
% \Acceptor{Quarta Castor}{Professor and Graduate Officer, Department of Terraforming}{}


% Date degree will be issued: \DegreeDate{Month}{year}
% Valid degree months at MIT are February, May, June, or September

\DegreeDate{June}{1876}


% Date that final thesis is submitted to department
\ThesisDate{May 18, 1876}


% If at least one thesis reader is included, a committee members page will be automatically generated.
% The MIT Libraries do not set a format for the committee members page, and if you prefer to use your own format,
% omit ALL readers here and insert your page just before the abstract.

\Reader{Marcus Gavius Apicius}{Professor of Cooking Arts}{Department of Food Science}
\Reader{Marie-Antoine Carême}{Professor of Haute Cuisine}{Department of Food Science}
\Reader{Miles Gloriosus}{Professor of Personal Pronouns}{Department of Rhetoric}

%% Are your thesis readers running off the page?  You can move the page heading upward by uncommenting this command:
% \CMPShortenTop


%%%%%%  Choose whether to have a CREATIVE COMMONS License  %%%%%%%%%%%%%%%%%%%%%%%%%%%%%%%%%%%%%%
%
% If you are using a cc license, uncomment the following line and insert details of your cc license here.
%
% \CClicense{CC BY-NC-ND 4.0}{https://creativecommons.org/licenses/by-nc-nd/4.0/}
%


%%%%%%%  Solutions for overflowing titlepage  %%%%%%%%%%%%%%%%%%%%%%%%%%%%%%%%%%%%%%%%%%%%%%%%%%%

% If your title page is overflowing (from too many names, degrees, etc.):
%
% (a) you can reduce the 12pt and 18pt skips between various blocks to 6pt with this command:
%
% \Tighten
%
% (b)  you can scale down the Signature block at the bottom with this command:
%
% \SignatureBlockSize{\small}  % or this one \SignatureBlockSize{\footnotesize}
%
% (c) you can put the acceptor name and title onto two lines, rather than three like this:
%
% \Acceptor{Tertius Castor}{Professor and Graduate Officer, Department of Research}{}
%
% (d) you can change the font size of the author name[s] with
%
%	\AuthorNameSize{\normalsize}
%
% (e) and you can omit any previous degrees from the title page, instead mentioning them in the biographical sketch

% Also, if you prefer to keep the text toward the top of the page with most white space at the bottom, you
% can use this command to squash all of the vertical glue (stretchy space) with this command:
%
% \Squash 
%
% This command is only useful when the text has not already reach the bottom of the page, since the glue gets squashed automatically
% when the page is very full.


%%%%%%%%%%%%%%%%%%%%%%%%%%%%%%%%%%%%%%%%%%%%%%%%%%%%%%%%%%%%%%%%%%%%%%%%%%%%%%%%%%%%%%%%%%%%%%%%%

%%% Make titlepage
\maketitle


%%%%%%%%  Content that you need to write follows!  %%%%%%%%%%%%%%%%%%%%%%%%%%%%%%%%%%%%%%%%%%%%%%

% \includeonly{acknowledgments,biography,chapter1,chapter2,...,appendixa,...} 
%   for usage of includeonly, see https://latexref.xyz/_005cinclude-_0026-_005cincludeonly.html


%%% Frontmatter (write this material in the mentioned files)  %%%%%%%%%%%%%%%%%%%%%%%%%%%%%%%%%%%

% If you are not using the automatic committee members page (by completing at least one \Reader macro), 
% you can insert your own page here. Create a file committee_members.tex and uncomment the next line:
% \include{committee_members}% This page is optional. 

% The abstract environment creates all the required headings and footers. 
% You only need to the text of the abstract in the file abstract.tex
\begin{abstract}
	\input{abstract.tex}% using \input, rather than \include, because we're inside an environment
\end{abstract}

\include{acknowledgments}% acknowledgments.tex (.tex extension is presumed by \include) 

\include{biography}% biography.tex (optional, see https://libraries.mit.edu/distinctive-collections/thesis-specs/#format)

%%% Table of contents and lists of stuff (delete unused lists, i.e., if no tables or figures) %%%%%

\tableofcontents
\listoffigures
\listoftables

%%% Chapters of thesis  %%%%%%%%%%%%%%%%%%%%%%%%%%%%%%%%%%%%%%%%%%%%%%%%%%%%%%%%%%%%%%%%%%%%%%%%%%%

%% If you want to use "double spacing", start here...

% From mitthesis package
% Version: 1.24, 2026/08/21
% Documentation: https://ctan.org/pkg/mitthesis


\chapter{Introduction}

\lipsum[1-2] Postremo aliquos futuros suspicor, qui me ad alias litteras vocent, genus hoc scribendi, etsi sit elegans, personae tamen et dignitatis esse negent~\cite{DKE1969,ww1920,kirk2288a,churchill1948,gibbs1863}.

\newcommand*{\Jpsi}{\ifpdftex\mathord{\bm{J}/\bm{\psi}}\else\mathord{\symbfit{J/\psi}}\fi} % works with either pdftex or lualatex

\section[A section discussing the first issue: \(J/\psi\)]{A section discussing the first issue: \( \Jpsi \) }

We begin with some ideas from the literature \cite{Fong2015,sharpe1}. 
\begin{equation}
\frac{\partial}{\partial t}\left[\rho\bigl(e + \lvert\vec{u}\rvert^2\big/2\bigr)\right]  + \nabla\cdot\left[\rho\bigl(h + \lvert\vec{u}\rvert^2\big/2 \bigr)\vec{u}\right]
 ={}-\nabla \cdot \vec{q} +  \rho \vec{u}\cdot\vec{g}+ \frac{\partial}{\partial x_j}\bigl(d_{ji}u_i\bigr)
\end{equation}
 \lipsum[3]

\lipsum[4] And more citations~\cite{sharpe1,GSL}.  Then we write some more and include our citations~\cite{Swaminathan2017IDABRO,dlmf,amsmath}. The configuration is shown in Figs.\ \ref{fig:golden} and~\ref{fig:golden1}.

%%%%%%%%%%%%%%%%%  begin figure  %%%%%%%%%%%%%%%%%%%%%%%%%%%
\begin{figure}[t]
% sample images are from mwe package, but should be found by latex in the tex tree w/o loading that package
\begin{subfigure}[t]{0.495\textwidth}
{\centering{\includegraphics[alt={A large letter C inside a rectangle},width=0.99\textwidth]{example-image-c.jpg}}}%
\subcaption{\label{fig:golden} Postremo aliquos futuros suspicor, qui me ad alias litteras vocent, genus hoc scribendi.}
\end{subfigure}
%%%%%%%% don't leave a break here
\begin{subfigure}[t]{0.495\textwidth}
{\centering{\includegraphics[alt={A large letter C inside a rectangle},width=0.99\textwidth]{example-image-c.jpg}}}%
\subcaption{\label{fig:golden1} Second subfigure.}%
\end{subfigure}%
\caption{A figure with two subfigures.\label{fig:4}}\label{fig:golden2}
\end{figure}
%%%%%%%%%%%%%%%%%%%  end figure  %%%%%%%%%%%%%%%%%%%%%%%%%%%%%

\lipsum[4]

%% note use of \ref* here to avoid placing a nested link in the table of contents
\subsection[Subsection~eqn.~(\ref*{eqn:WT1})]{Subsection~eqn.~\eqref{eqn:WT1}}
\lipsum[5-6]

\subsubsection{A subsubsection}
\lipsum[7]

\newcommand*{\boldA}{\ifpdftex\mathbf{A}\else\symbfup{A}\fi}% to get the right symbol with either pdftex or unicode-math

\begin{equation}\label{eqn:WT1}
L(\boldA) = \begin{pmatrix}
\dfrac\varphi{(\varphi_1,\varepsilon_1)}			& 0 												 & \ldots 									& \ldots & \ldots& 0 \\[4\jot]
\dfrac{\varphi k_{2,1}}{(\varphi_2,\varepsilon_1)}	& \dfrac\varphi{(\varphi_2,\varepsilon_2)}			 & 0 										& \ldots & \ldots& 0 \\[4\jot]
\dfrac{\varphi k_{3,1}}{(\varphi_3,\varepsilon_1)}	& \dfrac{\varphi k_{3,2}}{(\varphi_3,\varepsilon_2)} & \dfrac\varphi{(\varphi_3,\varepsilon_3)}	& 0 	 & \ldots& 0 \\[4\jot]
\vdots 												& \vdots 											 & \mbox{ } & \ddots & \mbox{ } & \vdots \\[\jot]
\dfrac{\varphi k_{n-1, 1}}{(\varphi_{n-1},\varepsilon_1)}		& \dfrac{\varphi k_{n-1, 2}}{(\varphi_{n-1},\varepsilon_2)} & \ldots & 
\dfrac{\varphi k_{n-1,n-2}}{(\varphi_{n-1},\varepsilon_{n-2})}	& \dfrac{\varphi}{(\varphi_{n-1},\varepsilon_{n-1})} 		& 0 \\[4\jot]
\dfrac{\varphi k_{n,1}}{(\varphi_n,\varepsilon_1)}				& \dfrac{\varphi k_{n,2}}{(\varphi_n,\varepsilon_2)}		& \ldots & \ldots	&
\dfrac{\varphi k_{n,n-1}}{(\varphi_n,\varepsilon_{n-1})} 		& \dfrac{\varphi}{(\varphi_n,\varepsilon_n)}
\end{pmatrix}
\end{equation}

%% This version looks nicer, but it does not produce proper html under current tagpdf (2025/10/31)
%
%\begin{equation}\label{eqn:WT1} 
%L(\boldA) = \begin{pmatrix}
%\dfrac\varphi{(\varphi_1,\varepsilon_1)}			& 0 												 & \hdotsfor{3} 							& 0 \\[4\jot]
%\dfrac{\varphi k_{2,1}}{(\varphi_2,\varepsilon_1)}	& \dfrac\varphi{(\varphi_2,\varepsilon_2)}			 & 0 										& \hdotsfor{2} & 0 \\[4\jot]
%\dfrac{\varphi k_{3,1}}{(\varphi_3,\varepsilon_1)}	& \dfrac{\varphi k_{3,2}}{(\varphi_3,\varepsilon_2)} & \dfrac\varphi{(\varphi_3,\varepsilon_3)}	& 0 & \hdotsfor{1} & 0 \\[\jot]
%\vdots 												&  													 &  & \smash{\rotatebox{15}{$\ddots$}} &  & \vdots \\[\jot]
%\dfrac{\varphi k_{n-1, 1}}{(\varphi_{n-1},\varepsilon_1)}		& \dfrac{\varphi k_{n-1, 2}}{(\varphi_{n-1},\varepsilon_2)} & \cdots        & 
%	\dfrac{\varphi k_{n-1,n-2}}{(\varphi_{n-1},\varepsilon_{n-2})}	& \dfrac{\varphi}{(\varphi_{n-1},\varepsilon_{n-1})} 		& 0 \\[4\jot]
%\dfrac{\varphi k_{n,1}}{(\varphi_n,\varepsilon_1)}				& \dfrac{\varphi k_{n,2}}{(\varphi_n,\varepsilon_2)}		& \hdotsfor{2}	&
%	\dfrac{\varphi k_{n,n-1}}{(\varphi_n,\varepsilon_{n-1})} 		& \dfrac{\varphi}{(\varphi_n,\varepsilon_n)}
%\end{pmatrix}
%\end{equation}

%\begin{equation}
%\int_{\Gamma^h_s}\frac{\partial}{\partial n_x}\int_{\Gamma^a_\xi} \mkern-3mu b(\xi) \frac{\partial E(x|\xi)}{\partial n_\xi} \,dl_\xi \Bigg\vert_{x\to s}\,dl_s = 0 
%\end{equation}

\section{Description our paradigm}\label{ch1:theidea}

\lipsum[8] No dissertation is complete without footnotes.\footnote{First footnote. $a_h = F_m$ See section~\ref{sec:stratified-flow}.}\footnote{Another interesting detail.}\footnote{And another really important idea to have in mind~\cite{reynolds1958,clauser56,lienhard2020,johnson1980,johnson1965,mpl}.} 

\begin{figure}[t]
% sample image is from mwe package, but should be found by latex in the tex tree w/o loading that package
\centering\includegraphics[alt={A large letter B inside a rectangle},width=6.67cm]{example-image-b.jpg} 
\caption{Caption text~\cite{GSL}.}\label{example-image-b}
\end{figure}


\subsection{Conversion to a metaheuristic}

\lipsum[11-12] This concept is discussed further in section~\ref{sec:stratified-flow}, and Refs.~\cite{euler1740,fourier1822}.


\section{Other generalizations}

\lipsum[7] And another citation, so that our sources will be unambiguous~\cite{montijano2014}.

\subsection{The most general case}
\lipsum[13]  But, according to Lord Brouncker (1655),
\begin{equation}
\frac{4}{\pi} = 1+ \cfrac{1^2}{2+
                     \cfrac{3^2}{2+
                       \cfrac{5^2}{2+
                         \cfrac{7^2}{2+
                           \cfrac{9^2}{2+\dotsb
}}}}}
\end{equation}

\section{Baroclinic generation of vorticity}\label{sec:stratified-flow}

Substitution of the particle acceleration and application Stokes theorem leads to the \textit{Kelvin-Bjerknes circulation theorem}, for
$\rho \neq \textrm{fn}(p)$:
\begin{align}
\frac{d\Gamma}{dt} &{}= \frac{d}{dt} \int_{\mathcal{C}} \mathbf{u} \cdot d\mathbf{r}\\
				   &{}= \int_{\mathcal{C}} \frac{D\mathbf{u}}{Dt} \cdot d\mathbf{r} + \underbrace{\int_{\mathcal{C}} \mathbf{u}\cdot d\biggl( \frac{d\mathbf{r}}{dt}\biggr)}_{=\, 0} \\[-2pt]
                   &{}= \iint_{\mathcal{S}} \nabla \times \frac{D\mathbf{u}}{Dt}  \cdot d\mathbf{A}\\
                   &{}= \iint_{\mathcal{S}}  \nabla p \times \nabla \left( \frac{1}{\rho}\right) \cdot d\mathbf{A}
\end{align}

Baroclinic generation of vorticity accounts for the sea breeze and various other atmospheric currents in which temperature, rather than pressure, creates density gradients. Further, this phenomenon accounts for ocean currents in straits joining more and less saline seas, with surface currents flowing from the fresher to the saltier water and with bottom current going oppositely.

\paragraph{Weak shock wave.} A weak shock wave travelling in a perfect gas is described by the following Burgers equation derived from approximations due to G. I. Taylor (1910)
\begin{equation}
u_t + \left(a_0 + \tfrac12(\gamma + 1)u\right) u_x = \frac{\delta}{2}u_{xx}
\end{equation}
Here, $\gamma$ is the specific heat capacity ratio, $\delta$ is the diffusivity of sound, $u_1$ is the downstream speed, $u_2 < u_1$ is the upstream speed, and $a_0$ is the undisturbed sound speed.  The  travelling-wave, or Taylor shock, solution of this equation is
\begin{equation}
\frac{u - u_2}{u_1-u_2} = \frac12\left[ 1 -\tanh\left( \frac{\gamma+1}{4\delta}(u_1-u_2)(x-ct) \right) \right]
\end{equation}
where $c = a_0 + (\gamma+1)(u_1+u_2)/4$. This solution exists for a compressive wave, $u_1 > u_2$.


%%%%%%%%%%%%%%% begin table %%%%%%%%%%%%%%%%%% 
\begin{table}[t]
\caption{The error function and complementary error function\label{tab:1}}%
\addvspace{0.5em}% backward compatibility for a pre-2026 class file
\centering{%
\tagpdfsetup{table/header-rows={1}}
\begin{tabular*}{0.8\textwidth}{@{\hspace*{1.5em}}@{\extracolsep{\fill}}ccc!{\hspace*{3.em}}ccc@{\hspace*{1.5em}}}
\toprule
\multicolumn{1}{@{\hspace*{1.5em}}c}{$x$\rule{0pt}{8pt}} &
$\erf{x}$ &
\multicolumn{1}{c!{\hspace*{3.em}}}{$\erfc{x}$} &
$x$       &
$\erf{x}$ &
\multicolumn{1}{c@{\hspace*{1.5em}}}{$\erfc{x}$} \\ \midrule
0.00 & 0.00000 & 1.00000 & 1.10 & 0.88021 & 0.11980 \\
0.05 & 0.05637 & 0.94363 & 1.20 & 0.91031 & 0.08969 \\
0.10 & 0.11246 & 0.88754 & 1.30 & 0.93401 & 0.06599 \\
0.15 & 0.16800 & 0.83200 & 1.40 & 0.95229 & 0.04771 \\
0.20 & 0.22270 & 0.77730 & 1.50 & 0.96611 & 0.03389 \\
0.30 & 0.32863 & 0.67137 & 1.60 & 0.97635 & 0.02365 \\
0.40 & 0.42839 & 0.57161 & 1.70 & 0.98379 & 0.01621 \\
0.50 & 0.52050 & 0.47950 & 1.80 & 0.98909 & 0.01091 \\
0.60 & 0.60386 & 0.39614 & 1.82\makebox[0pt][l]{14} & 0.99000 & 0.01000 \\
0.70 & 0.67780 & 0.32220 & 1.90 & 0.99279 & 0.00721 \\
0.80 & 0.74210 & 0.25790 & 2.00 & 0.99532 & 0.00468 \\
0.90 & 0.79691 & 0.20309 & 2.50 & 0.99959 & 0.00041 \\
1.00 & 0.84270 & 0.15730 & 3.00 & 0.99998 & 0.00002 \\
\bottomrule
\end{tabular*}
}%
\end{table}
%%%%%%%%%%%%%%%% end table %%%%%%%%%%%%%%%%%%% 

%%%%%%%%%%%%%%% begin more complicated table %%%%%%%%%%%%%%%%%%%%%%%%%%%%%%%%%%%%

% note column set as a 3 cm wide paragraph with 1 em hanging indentation. See array package documentation.
% note numbers centered on "." (with d{3.3}) and on "," (with ,{3.3}).  See dcolumn package documentation.
\begin{table}[t]
\caption{Table with more complicated columns: $\Delta T_m = \delta\theta/\Theta_0$\label{table:2a}}%
\addvspace{0.5em}% backward compatibility for a pre-2026 class file
\centering{%
\tagpdfsetup{table/header-rows={1}}%
\begin{tabular}{!{\hspace*{0.5cm}} >{\raggedright\hangindent=1em} p{5cm} d{3.3} @{\hspace*{1cm}} d{3.3} !{\hspace*{0.5cm}}}
\toprule
Experiment & \multicolumn{1}{c@{\hspace*{1cm}}}{$u$ [m/s]} & \multicolumn{1}{c!{\hspace*{0.5cm}}}{$T$ [\textdegree C]} \\
\midrule
The first test we ran this morning   & 124.3     &   68.3   \\
The second test we ran this morning  &  82.50    &  103.46  \\
Our competitor's test                &  72.321   &  141.384 \\
\bottomrule
\end{tabular}
}
\end{table}

%%%%%%%%%%%%%%%% end table  %%%%%%%%%%%%%%%%%%%%%%%%%%%%%%%%%%%% 

%%%%%%%%%%%%%%% begin table %%%%%%%%%%%%%%%%%% 

\begin{table*}[t]
\caption{Brinkman number at which liquid dissipation makes $<5\%$ or $<1\%$ contribution to $\theta_{b}$ or $q_{w}$ from the Graetz solution at the points where $\theta_{b}$ equals 0.05, 0.02, or 0.01.\label{table:3a}}%
\addvspace{0.5em}% backward compatibility for a pre-2026 class file
\centering{%
\tagpdfsetup{table/header-rows={1,2}}%
\begin{tabular}{d{1.2}!{\hspace*{0.3em}}d{1.5}d{1.4}>{$}c<{$}>{$}c<{$}>{$}c<{$}>{$}c<{$}}
\toprule
&&& \multicolumn{4}{c}{$\vert\textrm{Br}\vert$} \\
\cmidrule(l{0.5em}){4-7}
\multicolumn{1}{c}{$\theta_{b}$} & \multicolumn{1}{!{\hspace*{0.3em}}c}{$x^{+}$} & \multicolumn{1}{c}{$q_{w}D_{h}/k\Delta T$} & 
		\textrm{5\% of $\theta_{b}$} & \textrm{1\%  of $\theta_{b}$} &  \textrm{5\% of $q_{w}$} & \textrm{1\% of $q_{w}$} 
\\
\midrule
0.05 & 0.09621 & 0.3770  & 3.646\times 10^{-3} & 7.292\times 10^{-4} & 1.571\times 10^{-3}& 3.142\times 10^{-4} \\
0.02 & 0.1266  & 0.1508  & 1.458\times 10^{-3} & 2.917\times 10^{-4} & 6.824\times 10^{-4}& 1.257\times 10^{-4} \\
0.01 & 0.1496  & 0.07541 & 7.292\times 10^{-4} & 1.458\times 10^{-4} & 3.142\times 10^{-4}& 6.284\times 10^{-5} \\
\bottomrule
\end{tabular}
}
\end{table*}

%%%%%%%%%%%%%%%%% end table  %%%%%%%%%%%%%%%%%%%%%%%%%%%%%%% 


\section{Summary}
\lipsum[14-16]

%%	Nomenclature list is optional
%
%	This environment takes four optional arguments:
%		[1] adjust space between symbol and definition
%		[2] name (heading) of the nomenclature list
%		[3] level - can be "section" or "chapter" depending on whether you
%			have one nomenclature list for whole thesis or one for each
%			chapter. 
%		[4] style. The default matches the style of [3], but you can 
%			instead choose [frontmatter] or [backmatter] if desired.
%
%	For a single-column nomenclature list, use \begin{nomenclature}
%	For a two-column nomenclature list, use \begin{nomenclature*} AND
%		**put \usepackage{multicol}** in your preamble (in the main .tex file)
%
	\newcommand*{\boldomega}{\ifpdftex\mathord{\bm{\omega}}\else\mathord{\symbfup{\omega}}\fi} % works with pdftex or lualatex
%
\begin{nomenclature*}[2em][Nomenclature for Chapter 1][section]
\EntryHeading{Roman letters}
\entry{$\mathcal{C}$}{material curve}
\entry{$\mathbf{r}$}{material position [m]}
\entry{$\mathbf{u}$}{velocity [m s$^{-1}$]}% this more cumbersome superscripting is [arguably..] better for html translation than $^{-1}$.
\EntryHeading{Greek letters}
\entry{$\Gamma$}{circulation [m$^2$ s$^{-1}$]}
\entry{$\rho$}{mass density [kg m$^{-3}$]}
\entry{$\boldomega$}{vorticity [s$^{-1}$]} 
\end{nomenclature*}
% .tex extension is presumed
% \include{chapter2}
% \include{chapter3}
% \include{chapter4}


%%% Appendices of thesis  %%%%%%%%%%%%%%%%%%%%%%%%%%%%%%%%%%%%%%%%%%%%%%%%%%%%%%%%%%%%%%%%%%%%%%%%

\appendix
\include{appendixa}%  longtable and sidewaystable examples
%% From mitthesis package
% Version: 1.03, 2026/08/21
% Documentation: https://ctan.org/pkg/mitthesis

%%%%%%%%%%%%%%%%%%%%%%%%%%%%%%%%%%%%%%%%%%%%%%%%%%%%%%%%%%%%%%%%%

%% NB: the mchem package is not compatible with pdf tagging 
%%     under TeX Live 2026 and earlier formats! Fatal errors may 
%%     result from combining these! 

%% NB: the listings package is not compatible with pdf tagging 
%%     under TeX Live 2026 and earlier formats! Fatal errors may 
%%     result from combining these! 

%%%%%%%%%%%%%%%%%%%%%%%%%%%%%%%%%%%%%%%%%%%%%%%%%%%%%%%%%%%%%%%%%

\chapter{Some additional \LaTeX{} packages}

\section{The \texttt{mhchem} package for chemical equations}

These examples of chemical formul\ae\ are copied directly from the documentation of the \texttt{mhchem} package, which was used to typeset them.
\begin{gather}
\ce{x Na(NH4)HPO4 ->[\Delta] (NaPO3)_x + x NH3 ^ + x H2O} \\[0.5em]
\ce{^234_90Th -> ^0_-1$\beta${} + ^234_91Pa} \\[0.5em]
\ce{SO4^2- + Ba^2+ -> BaSO4 v} \\[0.5em]
\ce{Zn^2+
<=>[+ 2OH-][+ 2H+]
$\underset{\textrm{amphoteric hydroxide}}{\ce{Zn(OH)2 v}}$
<=>[+ 2OH-][+ 2H+]
$\underset{\textrm{tetrahydroxozincate}}{\ce{[Zn(OH)4]^2-}}$
}
\end{gather}


\section{The \texttt{listings} package for computer code}

This example uses the \texttt{listings} package.

\bigskip

\lstdefinestyle{mystyle}{
    backgroundcolor=\color{CadetBlue!15!white},   
    commentstyle=\color{Red3},
    numberstyle=\tiny\color{gray},
    stringstyle=\color{Blue3},
    basicstyle=\small\ttfamily,
    breakatwhitespace=false,         
    breaklines=true,                 
    numbers=left,                    
    numbersep=5pt,                  
    showspaces=false,                
    showstringspaces=false,
    showtabs=false,                  
    tabsize=2
}%
\lstset{language=[5.3]Lua,style={mystyle}}%

\begin{lstlisting}
\begin{luacode*}
function print_rate(kappa,xMin,xMax,npoints,option)
     local c = 1-kappa*kappa
     local croot = (1-kappa*kappa)^(1/2)
     local logx = math.log(xMin)
     local psi = 0
     
     local xstep = (math.log(xMax)-math.log(xMin))/(npoints-1)
     
     arg0 = math.sqrt(xMin/c)
     psi0 = (1/c)*math.exp((kappa*arg0)^2)*(erfc(kappa*arg0)-erfc(arg0))
     
     if option~=[[]] then
  		 tex.sprint("\\addplot+["..option.."] coordinates{") 
  		 -- addplot+ for color cycle to work
     else
  		 tex.sprint("\\addplot+ coordinates{")
     end
     tex.sprint("("..xMin..","..psi0..")")
     
     for i=1, (npoints-1) do
  		 x = math.exp(logx + xstep)
  		 arg = math.sqrt(x/c)
  		 karg = kappa*arg
  		 if karg<5 then 
		 -- this break compensates for exp(karg^2), which multiplies the error in the erf approximation...
  		    logpsi = -math.log(croot) + karg^2 + math.log(erfc(karg)-erfc(arg))
  		    psi = math.exp(logpsi)
  		 else
  		    psi = (1/(karg) - 1/(2*(karg^3)) + 3/(4*(arg^5)) )/(1.77245385*croot)
  		    -- this is the large x asymptote of the reaction rate
  		 end
  		 logx = math.log(x)
  		 tex.sprint("("..x..","..psi..")")
     end
     tex.sprint("}")
end
\end{luacode*}
\end{lstlisting}
% not compatible with tagged pdf as of June 2026


%%% Bibliography (biblatex)  %%%%%%%%%%%%%%%%%%%%%%%%%%%%%%%%%%%%%%%%%%%%%%%%%%%%%%%%%%%%%%%%%%%%%%

\defbibheading{bibintoc}{\chapter*{#1}\addcontentsline{toc}{backmatter}{\refname}} 
% this sets the title of contents name for bibliography to \refname (= References)
% change "backmatter" to "chapter" if you prefer a bold face entry in the table of contents

\printbibliography[title={\refname},heading=bibintoc]

% biblatex also supports chapter-by-chapter bibliography, https://tex.stackexchange.com/a/296502/119566
% see the biblatex manual, section 3.14.3


\end{document} 
 